\chapter{Einleitung} \label{ch:Einleitung}

\section{Aufgabe und Motivation}
\label{sec:AuM}

Das Ziel dieses physikalischen Praktikumsversuchs (Versuch 243) ist es, einerseits das fundamentale Phänomen des thermischen Rauschens in elektrischen Leitern zu verstehen und andererseits die Boltzmann-Konstante $k_B$ mittels einer hochpräzisen elektrischen Messmethode zu bestimmen. Die Motivation für diesen Versuch liegt in der Erkenntnis, dass die Genauigkeit elektrischer Messungen letztlich durch das thermische Rauschen begrenzt wird. Die zufällige Bewegung von Ladungsträgern führt zu unvermeidbaren Spannungsfluktuationen, die als Rauschuntergrund jede Messung überlagern. Um Messungen in der Wissenschaft und Messtechnik zu optimieren, ist es essenziell, die Ursachen dieses Rauschens zu kennen und Maßnahmen zu ergreifen, um das Signal-zu-Rausch-Verhältnis zu maximieren. Dies kann beispielsweise durch die Verwendung rauscharmer Verstärker, Kühlung der Messobjekte oder die gezielte Begrenzung der Messbandbreite erreicht werden.

Interessanterweise wird das thermische Rauschen nicht nur als Störgröße betrachtet, sondern findet auch gezielte Anwendung, etwa in Rauschthermometern, die Temperaturbereiche von wenigen Millikelvin bis zu mehreren tausend Kelvin abdecken können. In diesem Versuch wird das Rauschen genutzt, um eine der fundamentalen Naturkonstanten, die Boltzmann-Konstante, experimentell zu ermitteln. Die Bestimmung erfolgt indirekt über die Messung der Rauschspannung an ohmschen Widerständen in Abhängigkeit von deren Widerstandswert und Temperatur. Da die zu messenden Spannungen im Mikrovolt-Bereich liegen, ist ein rauscharmer Verstärker mit hoher Verstärkung notwendig, dessen eigenes Eigenrauschen dabei sorgfältig berücksichtigt und vom Gesamtsignal abgezogen werden muss.

\section{Physikalische Grundlage}
\label{sec:PhyBasics}

Thermisches Rauschen, auch als Johnson-Nyquist-Rauschen bezeichnet, ist ein statistisches Phänomen, das in allen elektrischen Leitern auftritt, sofern diese eine Temperatur oberhalb des absoluten Nullpunkts besitzen. Die Ursache liegt in der thermischen Energie der Ladungsträger, die eine ungeordnete, zufällige Bewegung ausführen. Diese Bewegung wird oft als Brownsche Bewegung der Ladungsträger beschrieben. Selbst ohne Anlegung einer äußeren Spannung führt diese ungerichtete Bewegung zu einem statistisch variierenden elektrischen Potenzial im Leiter. Misst man die Spannung über einem Widerstand mit einem hochempfindlichen Oszilloskop, so beobachtet man eine fluktuierende Spannung $U_r(t)$, die statistisch um einen Mittelwert schwankt. Da die Bewegung der Ladungsträger ungerichtet ist, heben sich positive und negative Beiträge im zeitlichen Mittel auf, sodass der Erwartungswert der Rauschspannung verschwindet:

\eq{mean_zero}{\langle U_r \rangle = \lim_{t' \to \infty} \frac{1}{t'} \int_0^{t'} U_r(t) \, dt = 0.}

Um die Stärke des Rauschens dennoch zu quantifizieren, muss daher nicht der Mittelwert, sondern der quadratische Effektivwert (Root Mean Square, rms) herangezogen werden. Dieser entspricht der Wurzel aus dem zeitlichen Mittelwert des Quadrats der Spannung:

\eq{rms_def}{\sqrt{\langle U_r^2 \rangle} = \sqrt{\lim_{t' \to \infty} \frac{1}{t'} \int_0^{t'} U_r(t)^2 \, dt}.}

Bei der Untersuchung des Frequenzspektrums einer solchen thermischen Rauschquelle stellt man fest, dass alle Frequenzanteile bis in den Terahertz-Bereich in gleichem Maße vorhanden sind. Keine Frequenz wird gegenüber einer anderen bevorzugt. In Analogie zum weißen Licht, bei dem alle sichtbaren Frequenzen gleich stark vertreten sind, wird dieses Phänomen als weißes Rauschen bezeichnet. Harry Nyquist leitete 1927 eine Beziehung her, die den quadratischen Effektivwert der Rauschspannung $\langle U_r^2 \rangle$ eines ohmschen Widerstands $R$ bei der Temperatur $T$ in Abhängigkeit von der Messbandbreite $\Delta f$ beschreibt:

\eq{nyquist}{\langle U_r^2 \rangle = 4 k_B T R \Delta f.}

In dieser Gleichung stellt $k_B$ die gesuchte Boltzmann-Konstante dar. Die Bandbreite $\Delta f$ bezieht sich hierbei auf die Bandbreite der verwendeten Messelektronik (Verstärker, Filter, Messgerät). Da das Rauschen frequenzunabhängig ist, trägt jeder Frequenzbereich gleichermaßen zur gesamten Rauschleistung bei; eine größere Bandbreite führt somit zu einer höheren gemessenen Rauschspannung.

Da die Rauschspannung bei Zimmertemperatur und typischen Widerstandswerten (z. B. einige k$\Omega$) und Bandbreiten nur im Bereich von Mikrovolt liegt, ist eine direkte Messung ohne Verstärkung nicht möglich. Ein rauscharmer Verstärker mit einer Verstärkung $V$ (im Versuch ca. 1000-fach oder 60 dB) wird daher eingesetzt. Allerdings stellt der Verstärker selbst eine weitere Rauschquelle dar. Die Rauschbeiträge des Widerstands $\langle U_R^2 \rangle$ und des Verstärkers $\langle U_V^2 \rangle$ addieren sich quadratisch, da sie statistisch unabhängig voneinander sind und ihre Mittelwerte null sind:

\eq{total_noise}{\langle U_{R+V}^2 \rangle = \langle (U_R + U_V)^2 \rangle = \langle U_R^2 \rangle + \langle U_V^2 \rangle.}

Um die Boltzmann-Konstante bestimmen zu können, muss das Eigenrauschen des Verstärkers $\langle U_V^2 \rangle$ experimentell ermittelt werden, indem der Verstärkereingang kurzgeschlossen wird (Nullmessung), und später von der Gesamtmessung subtrahiert werden. Zudem wird zur präzisen Bestimmung der effektiven Bandbreite ein Bandfilter verwendet, das aus einem Hochpass- und einem Tiefpassfilter besteht. Der Hochpass unterdrückt niederfrequente Störungen wie das 50-Hz-Netzbrummen, während der Tiefpass hochfrequentes Rauschen filtert und eine scharfe Begrenzung der oberen Grenzfrequenz ermöglicht.

Der Frequenzgang $g(f)$ des gesamten Messsystems (Verstärker plus Bandfilter) ist definiert als das Verhältnis der Ausgangsspannung zur Eingangsspannung bei einer bestimmten Frequenz $f$:

\eq{freq_response}{g(f) = \left. \frac{U_{\text{aus}}(f)}{U_{\text{ein}}(f)} \right|_f.}

Die Ausgangsrauschleistung ergibt sich durch Integration des quadratischen Frequenzgangs über alle Frequenzen, multipliziert mit der Nyquist-Leistungsdichte. Unter Berücksichtigung des Verstärkerrauschens gilt für die gemessene Ausgangsspannung:

\eq{output_voltage}{\langle U_{\text{aus}}^2 \rangle = 4 k_B T R \int_0^\infty g(f)^2 \, df + \langle U_V^2 \rangle = 4 k_B T R B + \langle U_V^2 \rangle.}

Das Integral $B = \int_0^\infty g(f)^2 \, df$ wird als äquivalente Rauschbandbreite des Messsystems bezeichnet. Durch Umstellen dieser Gleichung lässt sich die Boltzmann-Konstante berechnen:

\eq{boltzmann_calc}{k_B = \frac{\langle U_{\text{aus}}^2 \rangle - \langle U_V^2 \rangle}{4 T R B}.}

\subsubsection*{Skizze des Versuchsaufbaus}
\label{sec:Skizze}

Der Versuchsaufbau besteht aus einer Rauschquelle, die aus verschiedenen ohmschen Widerständen in einem abgeschirmten Gehäuse besteht. Diese werden direkt an den Eingang eines rauscharmen Verstärkers angeschlossen. Der Ausgang des Verstärkers ist mit einem Bandfilter verbunden, welches die gewünschte Messbandbreite definiert und Störsignale filtert. Das gefilterte Signal wird schließlich an ein hochempfindliches Effektivwert-Multimeter (Agilent HP34401A) oder ein Oszilloskop geleitet, um die Rauschspannung zu messen. Zur Bestimmung des Frequenzgangs wird ein Funktionsgenerator verwendet, der ein Sinussignal liefert. Dieses Signal wird vor dem Verstärker durch ein Dämpfungsglied (Dämpfungsfaktor $D = 10^{-3}$) stark abgeschwächt, um eine Übersteuerung (Sättigung) des Verstärkers zu vermeiden. Das Ausgangssignal des Bandfilters wird hierbei mit einem Oszilloskop gemessen. Die gesamte Steuerung der Messgeräte sowie die Datenerfassung erfolgen über einen PC.

